\documentclass[10pt,a4paper]{amsart}
\usepackage{amsmath,amssymb,amsthm,graphicx}
% Retained minimal preamble for the Stacks Project chapter more-morphisms.tex (More on Morphisms).
% Verbatim excerpt of the chapter preamble: Xy-pic setup (native preamble lines 14-18) and the native
% theorem declarations with the mathematical macros (native preamble lines 154-217). Copied mechanically;
% nothing added, removed or redefined. The wrapper class and the amsmath/amssymb/amsthm/graphicx base
% packages are supplied by the assembler, so the class/externaldocument part of the native preamble is omitted.
% For commutative diagrams we use Xy-pic
\usepackage[all]{xy}

% We use 2cell for 2-commutative diagrams.
\xyoption{2cell}
% Theorem environments.
%
\theoremstyle{plain}
\newtheorem{theorem}[subsection]{Theorem}
\newtheorem{proposition}[subsection]{Proposition}
\newtheorem{lemma}[subsection]{Lemma}

\theoremstyle{definition}
\newtheorem{definition}[subsection]{Definition}
\newtheorem{example}[subsection]{Example}
\newtheorem{exercise}[subsection]{Exercise}
\newtheorem{situation}[subsection]{Situation}

\theoremstyle{remark}
\newtheorem{remark}[subsection]{Remark}
\newtheorem{remarks}[subsection]{Remarks}

\numberwithin{equation}{subsection}

% Macros
%
\def\lim{\mathop{\mathrm{lim}}\nolimits}
\def\colim{\mathop{\mathrm{colim}}\nolimits}
\def\Spec{\mathop{\mathrm{Spec}}}
\def\Hom{\mathop{\mathrm{Hom}}\nolimits}
\def\Ext{\mathop{\mathrm{Ext}}\nolimits}
\def\SheafHom{\mathop{\mathcal{H}\!\mathit{om}}\nolimits}
\def\SheafExt{\mathop{\mathcal{E}\!\mathit{xt}}\nolimits}
\def\Sch{\mathit{Sch}}
\def\Mor{\mathop{\mathrm{Mor}}\nolimits}
\def\Ob{\mathop{\mathrm{Ob}}\nolimits}
\def\Sh{\mathop{\mathit{Sh}}\nolimits}
\def\NL{\mathop{N\!L}\nolimits}
\def\CH{\mathop{\mathrm{CH}}\nolimits}
\def\proetale{{pro\text{-}\acute{e}tale}}
\def\etale{{\acute{e}tale}}
\def\QCoh{\mathit{QCoh}}
\def\Ker{\mathop{\mathrm{Ker}}}
\def\Im{\mathop{\mathrm{Im}}}
\def\Coker{\mathop{\mathrm{Coker}}}
\def\Coim{\mathop{\mathrm{Coim}}}

% Boxtimes
%
\DeclareMathSymbol{\boxtimes}{\mathbin}{AMSa}{"02}

%
% Macros for moduli stacks/spaces
%
\def\QCohstack{\mathcal{QC}\!\mathit{oh}}
\def\Cohstack{\mathcal{C}\!\mathit{oh}}
\def\Spacesstack{\mathcal{S}\!\mathit{paces}}
\def\Quotfunctor{\mathrm{Quot}}
\def\Hilbfunctor{\mathrm{Hilb}}
\def\Curvesstack{\mathcal{C}\!\mathit{urves}}
\def\Polarizedstack{\mathcal{P}\!\mathit{olarized}}
\def\Complexesstack{\mathcal{C}\!\mathit{omplexes}}
% \Pic is the operator that assigns to X its picard group, usage \Pic(X)
% \Picardstack_{X/B} denotes the Picard stack of X over B
% \Picardfunctor_{X/B} denotes the Picard functor of X over B
\def\Pic{\mathop{\mathrm{Pic}}\nolimits}
\def\Picardstack{\mathcal{P}\!\mathit{ic}}
\def\Picardfunctor{\mathrm{Pic}}
\def\Deformationcategory{\mathcal{D}\!\mathit{ef}}
\usepackage{cleveref}
\crefname{originallemma}{original lemma}{original lemmas}
\Crefname{originallemma}{Original lemma}{Original lemmas}
\makeatletter
\newlabel{algebra-lemma-integral-overring-surjective}{{algebra-lemma-integral-overring-surjective}{source}{}{sourcepointer-1}{}}
\newlabel{algebra-lemma-integral-overring-surjective@cref}{{[originallemma][0][]algebra-lemma-integral-overring-surjective}{[1][0][]source}{}{}{}}
\makeatother
\crefname{originallemma}{original lemma}{original lemmas}
\Crefname{originallemma}{Original lemma}{Original lemmas}
\makeatletter
\newlabel{algebra-lemma-limit-module-finite-presentation}{{algebra-lemma-limit-module-finite-presentation}{source}{}{sourcepointer-2}{}}
\newlabel{algebra-lemma-limit-module-finite-presentation@cref}{{[originallemma][0][]algebra-lemma-limit-module-finite-presentation}{[1][0][]source}{}{}{}}
\makeatother
\crefname{originallemma}{original lemma}{original lemmas}
\Crefname{originallemma}{Original lemma}{Original lemmas}
\makeatletter
\newlabel{algebra-proposition-ubiquity-nagata}{{algebra-proposition-ubiquity-nagata}{source}{}{sourcepointer-3}{}}
\newlabel{algebra-proposition-ubiquity-nagata@cref}{{[originallemma][0][]algebra-proposition-ubiquity-nagata}{[1][0][]source}{}{}{}}
\makeatother
\crefname{originallemma}{original lemma}{original lemmas}
\Crefname{originallemma}{Original lemma}{Original lemmas}
\makeatletter
\newlabel{lemma-nowhere-dense-generic-fibre}{{lemma-nowhere-dense-generic-fibre}{source}{}{sourcepointer-4}{}}
\newlabel{lemma-nowhere-dense-generic-fibre@cref}{{[originallemma][0][]lemma-nowhere-dense-generic-fibre}{[1][0][]source}{}{}{}}
\makeatother
\crefname{originallemma}{original lemma}{original lemmas}
\Crefname{originallemma}{Original lemma}{Original lemmas}
\makeatletter
\newlabel{limits-lemma-descend-dimension-d}{{limits-lemma-descend-dimension-d}{source}{}{sourcepointer-5}{}}
\newlabel{limits-lemma-descend-dimension-d@cref}{{[originallemma][0][]limits-lemma-descend-dimension-d}{[1][0][]source}{}{}{}}
\makeatother
\crefname{originallemma}{original lemma}{original lemmas}
\Crefname{originallemma}{Original lemma}{Original lemmas}
\makeatletter
\newlabel{limits-lemma-limit-dimension}{{limits-lemma-limit-dimension}{source}{}{sourcepointer-6}{}}
\newlabel{limits-lemma-limit-dimension@cref}{{[originallemma][0][]limits-lemma-limit-dimension}{[1][0][]source}{}{}{}}
\makeatother
\crefname{originallemma}{original lemma}{original lemmas}
\Crefname{originallemma}{Original lemma}{Original lemmas}
\makeatletter
\newlabel{morphisms-lemma-dimension-fibre-after-base-change}{{morphisms-lemma-dimension-fibre-after-base-change}{source}{}{sourcepointer-7}{}}
\newlabel{morphisms-lemma-dimension-fibre-after-base-change@cref}{{[originallemma][0][]morphisms-lemma-dimension-fibre-after-base-change}{[1][0][]source}{}{}{}}
\makeatother
\crefname{originallemma}{original lemma}{original lemmas}
\Crefname{originallemma}{Original lemma}{Original lemmas}
\makeatletter
\newlabel{morphisms-lemma-dimension-fibre-at-a-point-additive}{{morphisms-lemma-dimension-fibre-at-a-point-additive}{source}{}{sourcepointer-8}{}}
\newlabel{morphisms-lemma-dimension-fibre-at-a-point-additive@cref}{{[originallemma][0][]morphisms-lemma-dimension-fibre-at-a-point-additive}{[1][0][]source}{}{}{}}
\makeatother
\crefname{originallemma}{original lemma}{original lemmas}
\Crefname{originallemma}{Original lemma}{Original lemmas}
\makeatletter
\newlabel{morphisms-lemma-nagata-normalization}{{morphisms-lemma-nagata-normalization}{source}{}{sourcepointer-9}{}}
\newlabel{morphisms-lemma-nagata-normalization@cref}{{[originallemma][0][]morphisms-lemma-nagata-normalization}{[1][0][]source}{}{}{}}
\makeatother
\crefname{originallemma}{original lemma}{original lemmas}
\Crefname{originallemma}{Original lemma}{Original lemmas}
\makeatletter
\newlabel{morphisms-proposition-generic-flatness}{{morphisms-proposition-generic-flatness}{source}{}{sourcepointer-10}{}}
\newlabel{morphisms-proposition-generic-flatness@cref}{{[originallemma][0][]morphisms-proposition-generic-flatness}{[1][0][]source}{}{}{}}
\makeatother
\title{Fibre dimension is preserved under the substitution a to a\textasciicircum{}n + b for large n}
\author{The Stacks Project Authors}
\date{stacks-project:more-morphisms.tex:local-snapshot-20261011; GNU FDL 1.2 or later}
\begin{document}
\maketitle
\noindent\textit{Editorial note.} Statement span is lines 8595-8610; proof span is lines 8612-8777, the complete author proof of 166 lines.\par\medskip
\noindent\textit{Editorial note.} The proof contains one align-style display and several inline formulas; the retained preamble carries the shared macro definitions.\par\medskip
\noindent\textit{Editorial note.} The proof cites several labelled results, which are prerequisites and are supplied as non-load-bearing navigation stubs.\par\medskip
\noindent\textit{Editorial note.} The author's proof contains an informal one-word sentence beginning the final paragraph of the induction; it is retained verbatim.\par\medskip
\noindent\textit{Editorial note.} extra\_packages is empty on purpose: the retained preamble already loads xy with options and a second declaration clashes.\par\medskip
\noindent\textit{Allowed standard prerequisites (editorial).}\begin{itemize}
\item Explicit prerequisite boundary. Finite type and finite presentation ring maps, dimension of fibres of a morphism to the affine line, filtered colimits of rings and the behaviour of dimension along them, and the dimension bound for quotients by an ideal. The chain of reductions and the induction are retained.
\end{itemize}\medskip
\ifdefined\hypertarget\hypertarget{sourcepointer-1}{}\fi
\noindent\textit{Source reference.} Original lemma algebra-lemma-integral-overring-surjective: sources/stacks-more-morphisms/pinned-original-source.tex, source lines 8777--8777. Referenced by the retained statement or proof but not part of the collected spans; supplied only as a navigation pointer. The argument does not depend on its content.\par\medskip
\ifdefined\hypertarget\hypertarget{sourcepointer-2}{}\fi
\noindent\textit{Source reference.} Original lemma algebra-lemma-limit-module-finite-presentation: sources/stacks-more-morphisms/pinned-original-source.tex, source lines 8777--8777. Referenced by the retained statement or proof but not part of the collected spans; supplied only as a navigation pointer. The argument does not depend on its content.\par\medskip
\ifdefined\hypertarget\hypertarget{sourcepointer-3}{}\fi
\noindent\textit{Source reference.} Original lemma algebra-proposition-ubiquity-nagata: sources/stacks-more-morphisms/pinned-original-source.tex, source lines 8777--8777. Referenced by the retained statement or proof but not part of the collected spans; supplied only as a navigation pointer. The argument does not depend on its content.\par\medskip
\ifdefined\hypertarget\hypertarget{sourcepointer-4}{}\fi
\noindent\textit{Source reference.} Original lemma lemma-nowhere-dense-generic-fibre: sources/stacks-more-morphisms/pinned-original-source.tex, source lines 6524--6530. Referenced by the retained statement or proof but not part of the collected spans; supplied only as a navigation pointer. The argument does not depend on its content.\par\medskip
\ifdefined\hypertarget\hypertarget{sourcepointer-5}{}\fi
\noindent\textit{Source reference.} Original lemma limits-lemma-descend-dimension-d: sources/stacks-more-morphisms/pinned-original-source.tex, source lines 8777--8777. Referenced by the retained statement or proof but not part of the collected spans; supplied only as a navigation pointer. The argument does not depend on its content.\par\medskip
\ifdefined\hypertarget\hypertarget{sourcepointer-6}{}\fi
\noindent\textit{Source reference.} Original lemma limits-lemma-limit-dimension: sources/stacks-more-morphisms/pinned-original-source.tex, source lines 8777--8777. Referenced by the retained statement or proof but not part of the collected spans; supplied only as a navigation pointer. The argument does not depend on its content.\par\medskip
\ifdefined\hypertarget\hypertarget{sourcepointer-7}{}\fi
\noindent\textit{Source reference.} Original lemma morphisms-lemma-dimension-fibre-after-base-change: sources/stacks-more-morphisms/pinned-original-source.tex, source lines 8777--8777. Referenced by the retained statement or proof but not part of the collected spans; supplied only as a navigation pointer. The argument does not depend on its content.\par\medskip
\ifdefined\hypertarget\hypertarget{sourcepointer-8}{}\fi
\noindent\textit{Source reference.} Original lemma morphisms-lemma-dimension-fibre-at-a-point-additive: sources/stacks-more-morphisms/pinned-original-source.tex, source lines 8777--8777. Referenced by the retained statement or proof but not part of the collected spans; supplied only as a navigation pointer. The argument does not depend on its content.\par\medskip
\ifdefined\hypertarget\hypertarget{sourcepointer-9}{}\fi
\noindent\textit{Source reference.} Original lemma morphisms-lemma-nagata-normalization: sources/stacks-more-morphisms/pinned-original-source.tex, source lines 8777--8777. Referenced by the retained statement or proof but not part of the collected spans; supplied only as a navigation pointer. The argument does not depend on its content.\par\medskip
\ifdefined\hypertarget\hypertarget{sourcepointer-10}{}\fi
\noindent\textit{Source reference.} Original lemma morphisms-proposition-generic-flatness: sources/stacks-more-morphisms/pinned-original-source.tex, source lines 8777--8777. Referenced by the retained statement or proof but not part of the collected spans; supplied only as a navigation pointer. The argument does not depend on its content.\par\medskip
\begin{lemma}
\label{lemma-sum-is-ok}
Let $R \to S$ be a finite type ring map. Let $d \geq 0$. Let $a, b \in S$.
Assume that the fibres of
$$
f_a : \Spec(S) \longrightarrow \mathbf{A}^1_R
$$
given by the $R$-algebra map $R[x] \to S$ sending $x$ to $a$
have dimension $\leq d$. Then there exists an $n_0$ such that for
$n \geq n_0$ the fibres of
$$
f_{a^n + b} : \Spec(S) \longrightarrow \mathbf{A}^1_R
$$
given by the $R$-algebra map $R[x] \to S$ sending $x$ to $a^n + b$
have dimension $\leq d$.
\end{lemma}

\begin{proof}
In this paragraph we reduce to the case where $R \to S$ is of finite
presentation. Namely, write $S = R[A, B, x_1, \ldots, x_n]/J$ for some ideal
$J \subset R[x_1, \ldots, x_n]$ where $A$ and $B$ map to $a$ and $b$ in $S$.
Then $J$ is the union of its finitely generated ideals $J_\lambda \subset J$.
Set $S_\lambda = R[A, B, x_1, \ldots, x_n]/J_\lambda$ and
denote $a_\lambda, b_\lambda \in S_\lambda$ the images of $A$ and $B$.
Then for some $\lambda$ the fibres of
$$
f_{a_\lambda} : \Spec(S_\lambda) \longrightarrow \mathbf{A}^1_R
$$
have dimension $\leq d$, see Limits, Lemma \ref{limits-lemma-limit-dimension}.
Fix such a $\lambda$. If we can find $n_0$ which works for $R \to S_\lambda$,
$a_\lambda$, $b_\lambda$, then $n_0$ works for $R \to S$. Namely, the fibres
of $f_{a_\lambda^n + b_\lambda} : \Spec(S_\lambda) \to \mathbf{A}^1_R$
contain the fibres of $f_{a^n + b} : \Spec(S) \to \mathbf{A}^1_R$.
This reduces us to the case discussed in the next paragraph.

\medskip\noindent
Assume $R \to S$ is of finite presentation. In this paragraph we reduce
to the case where $R$ is of finite type over $\mathbf{Z}$. By
Algebra, Lemma \ref{algebra-lemma-limit-module-finite-presentation}
we can find a directed set $\Lambda$ and a system of ring maps
$R_\lambda \to S_\lambda$ over $\Lambda$ whose colimit is
$R \to S$ such that $S_\mu = S_\lambda \otimes_{R_\lambda} R_\mu$
for $\mu \geq \lambda$ and such that each $R_\lambda$ and $S_\lambda$
is of finite type over $\mathbf{Z}$. Choose $\lambda_0 \in \Lambda$ and
elements $a_{\lambda_0}, b_{\lambda_0} \in S_{\lambda_0}$ mapping to
$a, b \in S$. For $\lambda \geq \lambda_0$ denote
$a_\lambda, b_\lambda \in S_\lambda$
the image of $a_{\lambda_0}, b_{\lambda_0}$.
Then for $\lambda \geq \lambda_0$ large enough the fibres of
$$
f_{a_\lambda} : 
\Spec(S_\lambda) \longrightarrow \mathbf{A}^1_{R_\lambda}
$$
have dimension $\leq d$, see
Limits, Lemma \ref{limits-lemma-descend-dimension-d}.
Fix such a $\lambda$. If we can find $n_0$ which works for
$R_\lambda \to S_\lambda$, $a_\lambda$, $b_\lambda$, then $n_0$
works for $R \to S$. Namely, any fibre of
$f_{a^n + b} : \Spec(S) \to \mathbf{A}^1_R$
has the same dimension as a fibre of
$f_{a_\lambda^n + b_\lambda} : \Spec(S_\lambda) \to \mathbf{A}^1_{R_\lambda}$
by Morphisms, Lemma \ref{morphisms-lemma-dimension-fibre-after-base-change}.
This reduces us the the case discussed in the next paragraph.

\medskip\noindent
Assume $R$ and $S$ are of finite type over $\mathbf{Z}$. In particular
the dimension of $R$ is finite, and we may use induction on $\dim(R)$.
Thus we may assume the result holds for all situations with
$R' \to S'$, $a$, $b$ as in the lemma with $R'$ and $S'$ of finite type
over $\mathbf{Z}$ but with $\dim(R') < \dim(R)$.

\medskip\noindent
Since the statement is about the topology of the spectrum of $S$
we may assume $S$ is reduced. Let $S^\nu$ be the normalization of $S$.
Then $S \subset S^\nu$ is a finite extension as $S$ is excellent, see
Algebra, Proposition \ref{algebra-proposition-ubiquity-nagata}
and Morphisms, Lemma \ref{morphisms-lemma-nagata-normalization}.
Thus $\Spec(S^\nu) \to \Spec(S)$ is surjective and finite
(Algebra, Lemma \ref{algebra-lemma-integral-overring-surjective}).
It follows that if the result holds for $R \to S^\nu$ and the
images of $a$, $b$ in $S^\nu$, then the result holds for $R \to S$, $a$, $b$.
(Small detail omitted.) This reduces us to the case discussed in the
next paragraph.

\medskip\noindent
Assume $R$ and $S$ are of finite type over $\mathbf{Z}$ and
$S$ normal. Then $S = S_1 \times \ldots \times S_r$ for some
normal domains $S_i$. If the result holds for each $R \to S_i$
and the images of $a$, $b$ in $S_i$, then the result holds for
$R \to S$, $a$, $b$. (Small detail omitted.)
This reduces us to the case discussed in the next paragraph.

\medskip\noindent
Assume $R$ and $S$ are of finite type over $\mathbf{Z}$ and
$S$ a normal domain. We may replace $R$ by the image of $R$ in $S$
(this does not increase the dimension of $R$).
This reduces us to the case discussed in the next paragraph.

\medskip\noindent
Assume $R \subset S$ are of finite type over $\mathbf{Z}$ and
$S$ a normal domain. Consider the morphism
$$
f_a : \Spec(S) \to \mathbf{A}^1_R
$$
The assumption tells us that $f_a$ has fibres of dimension $\leq d$.
Hence the fibres of $f : \Spec(S) \to \Spec(R)$ have dimension $\leq d + 1$
(Morphisms, Lemma \ref{morphisms-lemma-dimension-fibre-at-a-point-additive}).
Consider the morphism of integral schemes
$$
\phi : \Spec(S) \to \mathbf{A}^2_R = \Spec(R[x, y])
$$
corresponding to the $R$-algebra map $R[x, y] \to S$ sending $x$ to $a$
and $y$ to $b$. There are two cases to consider
\begin{enumerate}
\item $\phi$ is dominant, and
\item $\phi$ is not dominant.
\end{enumerate}
We claim that in both cases there exists an integer $n_0$
and a nonempty open $V \subset \Spec(R)$ such that for $n \geq n_0$
the fibres of $f_{a^n + b}$ at points $q \in \mathbf{A}^1_V$
have dimension $\leq d$.

\medskip\noindent
Proof of the claim in case (1). We have $f_{a^n + b} = \pi_n \circ \phi$
where
$$
\pi_n :  \mathbf{A}^2_R \to \mathbf{A}^1_R
$$
is the flat morphism corresponding to the $R$-algebra map $R[x] \to R[x, y]$
sending $x$ to $x^n + y$. Since $\phi$ is dominant there is a dense open
$U \subset \Spec(S)$ such that $\phi|_U : U \to \mathbf{A}^2_R$ is flat
(this follows for example from generic flatness, see
Morphisms, Proposition \ref{morphisms-proposition-generic-flatness}).
Then the composition
$$
f_{a^n + b}|_U :
U \xrightarrow{\phi|_U}
\mathbf{A}^2_R
\xrightarrow{\pi_n}
\mathbf{A}^1_R
$$
is flat as well. Hence the fibres of this morphism have at least
codimension $1$ in the fibres of $f|_U : U \to \Spec(R)$ by
Morphisms, Lemma \ref{morphisms-lemma-dimension-fibre-at-a-point-additive}.
In other words, the fibres of $f_{a^n + b}|_U$ have dimension $\leq d$.
On the other hand, since $U$ is dense in $\Spec(S)$, we can find a nonempty
open $V \subset \Spec(R)$ such that $U \cap f^{-1}(p) \subset f^{-1}(p)$
is dense for all $p \in V$ (see for example Lemma
\ref{lemma-nowhere-dense-generic-fibre}).
Thus $\dim(f^{-1}(p) \setminus U \cap f^{-1}(p)) \leq d$ and
we conclude that our claim is true (as any fibres of
$f_{a^n + b} : \Spec(S) \to \mathbf{A}^1_R$
is contained in a fibre of $f$).

\medskip\noindent
Case (2). In this case we can find a nonzero $g = \sum c_{ij} x^i y^j$
in $R[x, y]$ such that $\Im(\phi) \subset V(g)$. In fact, we may assume
$g$ is irreducible over $\text{Frac}(R)$. If $g \in R[x]$, say with leading
coefficient $c$, then over $V = D(c) \subset \Spec(R)$ the fibres of
$f$ already have dimension $\leq d$ (because the image of $f_a$ is contained
in $V(g) \subset \mathbf{A}^1_R$ which has finite fibres over $V$). Hence
we may assume $g$ is not contained in $R[x]$.
Let $s \geq 1$ be the degree of $g$ as a polynomial in $y$ and
let $t$ be the degree of $\sum c_{is} x^i$ as a polynomial in $x$.
Then $c_{ts}$ is nonzero and
$$
g(x, -x^n) = (-1)^s c_{ts} x^{t + sn} + l.o.t.
$$
provided that $n$ is bigger than the degree of $g$ as a polynomial in $x$
(small detail omitted). For such $n$ the polynomial $g(x, -x^n)$ 
is a nonzero polynomial in $x$ and maps to a nonzero polynomial in
$\kappa(\mathfrak p)[x]$ for all $\mathfrak p \subset R$,
$c_{st} \not \in \mathfrak p$. We conclude that our claim is
true for $V$ equal to the principal open $D(c_{ts})$ of $\Spec(R)$.

\medskip\noindent
OK, and now we can use induction on $\dim(R)$. Namely, let
$I \subset R$ be an ideal such that $V(I) = \Spec(R) \setminus V$.
Observe that $\dim(R/I) < \dim(R)$ as $R$ is a domain.
Let $n'_0$ be the integer we have by induction on $\dim(R)$
for $R/I \to S/IS$ and the images of $a$ and $b$ in $S/IS$.
Then $\max(n_0, n'_0)$ works.
\end{proof}

\end{document}
